% titel: Blatt 7, A2: Zug und Dehnung
% skript: Physik I, Kapitel 7.1--7.2, Seiten 122--123

Zieht eine Kraft an einem Werkstück, werden innerhalb des Materials Spannung
und Längenänderung erzeugt. Untersucht wird, wann das Werkstück bricht oder
sich dauerhaft verformt. Zwei Grössen genügen, um die Belastung zu fassen:
Spannung (Last, bezogen auf den Querschnitt) und Dehnung (relative
Längenänderung).

\begin{werkzeugkasten}[Spannung und Dehnung]
Eine über den Querschnitt gleichmässig verteilte Zugkraft $F$ auf der Fläche
$A$ erzeugt eine Normalspannung; die relative Längenänderung heisst Dehnung:
\begin{equation}
  \sigma=\frac{F}{A}, \qquad
  \varepsilon=\frac{\Delta L}{L_0}.
\end{equation}
$\varepsilon$ ist dimensionslos. Die Spannung hat die Einheit
$\si{\pascal}=\si{\newton\per\metre\squared}$. So wird aus einer Kraft und
einer Geometrie eine materialunabhängige Belastungsgrösse.

\textit{Skript:} Kap. 7.1, S. 122, Gl. (7.1) und (7.3).
\end{werkzeugkasten}

\begin{werkzeugkasten}[Hooke'sches Gesetz]
Im linear-elastischen Bereich sind Spannung und Dehnung proportional:
\begin{equation}
  \sigma=E\varepsilon .
\end{equation}
Der Elastizitätsmodul $E$ (gleiche Einheit wie Spannung) beschreibt die
Steifigkeit, nicht die Festigkeit. Diese Beziehung gilt nur unterhalb des
Geltungsbereichs; für das Bruchverhalten zählt die Grenze von $E$.

\textit{Skript:} Kap. 7.2, S. 122--123, Gl. (7.4), und Kap. 7, S. 121.
\end{werkzeugkasten}

\begin{achtung}[Darauf achten]
\begin{itemize}
  \item Spannung und Dehnung auf denselben Querschnitt bzw. dieselbe
        Bezugslänge beziehen.
  \item Einheiten einheitlich umrechnen: $\qty{1}{\newton\per\milli\metre\squared}
        =\qty{1e6}{\pascal}$.
\end{itemize}
\end{achtung}
